\section{Introduction}

\subsection{Motivation}
Objective Structured Clinical Examinations (OSCEs) are widely used in medical
education to assess clinical and communication skills. In this setting, the
quality of examiner feedback is central to learning. However, feedback may vary
substantially between examiners and can be vague, overly judgmental,
insufficiently specific, or difficult for students to act upon. Feedback that
does not refer to observable behavior can support encouragement, but it often
does not provide enough information for targeted improvement
\cite{hattie2007power,gigante2011getting,ramani2012twelve}.

This problem is particularly relevant in current examiner training. Examiners
are required to give structured two-minute feedback at every OSCE station, but
many have not received formal training in how to do this. The skill is not
trivial: high-quality feedback requires observation, prioritisation, concise
wording, balanced judgement, and an actionable recommendation within a short
time window. An examiner may know the expected clinical performance, but still
struggle to formulate feedback that is educationally useful.

For example, a statement such as ``good job'' or ``you need to be more
systematic'' may be well intentioned, but it does not tell the learner which
observed action was effective, which behavior should be changed, or how to
improve in the next OSCE station. Examiner Coach addresses this gap by treating
feedback itself as the object of analysis. The system does not grade the
student in the OSCE video; instead, it evaluates the quality of the examiner's
spoken feedback to that student. The project therefore sits at the intersection
of medical education, natural language processing, and human-computer
interaction.

\subsection{Problem Definition}
The problem addressed in this report is not simply that feedback can be poor. It
is that feedback delivery is a skill that requires practice, guidance, and
timely correction, while the current OSCE training context provides limited
support for individual rehearsal. The project is defined by four connected
challenges:
\begin{itemize}
    \item examiners are expected to provide structured two-minute feedback at
    every station, but many have received no formal feedback-delivery training;
    \item existing examiner training is often delivered in large-group
    workshops with limited individual practice and personalised correction;
    \item there is no scalable, on-demand tool that allows examiners to practise
    feedback delivery and receive immediate criterion-based evaluation;
    \item privacy constraints in medical education make commercial AI services
    unsuitable for many examiner-training scenarios, because patient-adjacent
    and institutional data must remain protected.
\end{itemize}

The resulting training gap is therefore clear: examiners are expected to give
structured, useful feedback, but lack a scalable and privacy-conscious way to
rehearse this skill before entering the examination setting.

\subsection{Aim and Contributions}
This project aims to develop a prototype AI system that supports examiner
feedback training by evaluating spoken feedback and providing formative
coaching. The main contributions are:
\begin{itemize}
    \item a structured six-criterion model for feedback-quality evaluation;
    \item a backend-centered modular RAG pipeline for evidence-grounded
    assessment;
    \item bilingual German and English evaluation and coaching support;
    \item a prototype comparison of RAG variants, including no-RAG, direct
    retrieval, and HyDE-based retrieval.
\end{itemize}

The report therefore asks how a modular RAG architecture can be adapted to a
medical education feedback-training scenario, and whether retrieval-based
variants provide clearer calibration than a prompt-only baseline. The intended
outcome is not a final clinical validation, but a technically inspectable
prototype that can support later expert review and user studies.
